\tzpanchor{0318}
\tzpentryheader{0318}{Constrained Dipole Field Configuration}

\begingroup\small
\begingroup\scriptsize
\setlength{\tabcolsep}{4pt}
\renewcommand{\arraystretch}{0.92}
\noindent\begin{tabularx}{\linewidth}{@{}p{0.24\linewidth}p{0.36\linewidth}X@{}}
\textbf{UUID} & \multicolumn{2}{l@{}}{\tzpcode{9560f8f8-7a6d-566a-bc41-fabdf9220a18}}\\
\textbf{PiID} & \tzpcode{1748815209209} & \textbf{TZPi}\quad \tzpcode{5}\quad \textbf{PiPE}\quad \tzpcode{325}\\
\textbf{RTEID} & \textbf{Poloidal $\theta$}\quad \tzpcode{1234.8646 rad} & \textbf{Toroidal $\phi$}\quad \tzpcode{02.0347 rad}\\
\textbf{TZPID} & \tzpcode{ID0318} & \textbf{Node Dependencies}\quad \tzpidref{0317}\\
\textbf{Registry} & \tzpcode{registry\_id\_0318} & \textbf{Scale}\quad transfer\\
\textbf{LOC Classification} & QA641 manifolds and topology & QC174.17 mathematical physics\\
\textbf{arXiv Classification} & Primary: math-ph & Cross-list: gr-qc, math.DG\\
\textbf{Primary Support} & Dipole Locations & Single Dipole Field\\
\textbf{Closure Support} & Total Constrained Field & \\
\end{tabularx}\par
\endgroup
\endgroup

\tzpsection{Canonical Statement}
Two magnetic dipoles of equal magnitude $\mu_{0}$, oriented along the $z$-axis with opposite directions, are located at

\tzpsection{Canonical Equation}
\[ r_{1}=(0,0,-d/2), \qquad r_{2}=(0,0,d/2). \] The field produced by a single dipole at position $r_{0}$ is \[ B_{\mathrm{dip}}(r;\mu,r_{0}) = \frac{\mu_{0}}{4\pi} \left[ \frac{3(r-r_{0})\bigl(\mu\cdot(r-r_{0})\bigr)}{\lvert r-r_{0}\rvert^{5}} - \frac{\mu}{\lvert r-r_{0}\rvert^{3}} \right]. \] The total internal field is \[ B_{\mathrm{total}}(r,t) = B_{\mathrm{dip}}(r;\mu_{1}(t),r_{1}) + B_{\mathrm{dip}}(r;\mu_{2}(t),r_{2}) + B_{\mathrm{constraint}}(r,t). \]

\[
\mathbf{r}_{+}=a\hat{\mathbf{z}},\qquad \mathbf{r}_{-}=-a\hat{\mathbf{z}}
\]
\[
\mathbf{B}_{\mathrm{dip}}(\mathbf{r};\mathbf{m})=\frac{\mu_{0}}{4\pi}\frac{3\hat{\mathbf{r}}(\mathbf{m}\cdot\hat{\mathbf{r}})-\mathbf{m}}{\lVert\mathbf{r}\rVert^{3}}
\]

\begin{minipage}[t]{0.49\linewidth}
\tzpsection{Canonical Symbol Key}
\begingroup
\small
\raggedright
\textbf{Source equation symbols.}\par
\begin{itemize}[leftmargin=1.35em,itemsep=1pt,topsep=1pt,parsep=0pt]
    \item $r_{1}$: source equation symbol.
    \item $r_{2}$: source equation symbol.
    \item $B_{\mathrm{dip}}$: source equation symbol.
    \item $\mu$: source equation symbol.
    \item $\mu_{0}$: vacuum permeability constant.
    \item $\pi$: source equation symbol.
    \item $r_{0}$: source equation symbol.
    \item $B_{\mathrm{total}}(r,t)$: source equation symbol.
    \item $t$: source equation symbol.
    \item $\mu_{1}(t)$: source equation symbol.
\end{itemize}
\endgroup
\end{minipage}\hfill
\begin{minipage}[t]{0.47\linewidth}
\tzpsection{Wolfram Form}
\begingroup
\scriptsize
\raggedright
\textbf{Source Wolfram model.}\par
\begin{Verbatim}[fontsize=\tiny,breaklines=true,breakanywhere=true]
(* TZPID ID0318 generated source Wolfram model *)
ClearAll[K0318, Cr0318, T, R, A, W, I, N];
eq0318$1 = HoldForm[TZPSequence[r1 == Tuple[0, 0, -d/2], r2 == Tuple[0, 0, d/2]]];
eq0318$2 = HoldForm[BDip[r; mu, r0] == (mu0/(4*Pi))[(3*(r - r0)*(mu*(r - r0)))/Abs[r - r0]^5 - mu/Abs[r - r0]^3]];
eq0318$3 = HoldForm[BTotal[r, t] == BDipTuple[r; mu1[t], r1] + BDipTuple[r; mu2[t], r2] + BConstraint[r, t]];
eq0318$4 = HoldForm[TZPSequence[r == ahat[z], r == -ahat[z]]];
eq0318$5 = HoldForm[BDip[r; m] == (mu0/(4*Pi))*((3*hat[r]*(m*hat[r]) - m)/lVertrrVert^3)];
eq0318$6 = HoldForm[B == BDip[r - r; m] + BDip[r - r; -m]];

K0318 = HoldForm[{eq0318$1, eq0318$2, eq0318$3, eq0318$4, eq0318$5, eq0318$6}];

N = "normalized TRAWIN closure object for Constrained Dipole Field Configuration";
I = "Total Constrained Field integration/comparison layer";
W = "Single Dipole Field wave or operator carrier";
A = "Dipole Locations admissible action/amplitude condition";
R = "Single Dipole Field rotation/curl preservation layer";
T = "Dipole Locations local source condition";

Cr0318[expr_] := HoldForm[N[I[W[A[R[T[expr]]]]]]];

Result0318 = Cr0318[K0318];
\end{Verbatim}
\endgroup
\end{minipage}

\par\vspace{0.45\baselineskip}
\noindent\begin{minipage}[t]{\linewidth}
\begingroup
\small
\raggedright
\textbf{TRAWIN Operator Key.}\par
\noindent\textbf{Time/localization operator:} $\mathsf{T}\!\left[\mathrm{local}\right]$ Dipole Locations local source condition.\par
\noindent\textbf{Rotation/curl operator:} $\mathsf{R}\!\left[\nabla\times\right]$ Single Dipole Field rotation/curl preservation layer.\par
\noindent\textbf{Action/amplitude operator:} $\mathsf{A}\!\left[\mathrm{admissible}\right]$ Dipole Locations admissible action/amplitude condition.\par
\noindent\textbf{Wave/d'Alembertian operator:} $\mathsf{W}\!\left[\Box\right]$ Single Dipole Field wave or operator carrier.\par
\noindent\textbf{Integration operator:} $\mathsf{I}\!\left[\int\right]$ Total Constrained Field integration/comparison layer.\par
\noindent\textbf{Normalization operator:} $\mathsf{N}\!\left[\mathrm{normalized}\right]$ normalized TRAWIN closure object for Constrained Dipole Field Configuration.\par
\endgroup
\end{minipage}

\clearpage
\noindent\textbf{[ID 0318] Constrained Dipole Field Configuration}\par
\vspace{0.35em}
\tzpsection{Derivation Layer}
\paragraph{Primary Equation.}
\begin{equation}
\Delta_T\phi
=
\nabla^{T}_{a}\nabla^{T\,a}\phi .
\end{equation}

\paragraph{Primary Derivation.}
For a scalar field $\phi$, the ordinary Laplacian is obtained by taking the divergence of the gradient.  
Replacing the ordinary connection with the Trawinistic connection gives:

\begin{equation}
\Delta_T\phi
=
\nabla^{T}_{a}
\left(
\nabla^{T\,a}\phi
\right).
\end{equation}

Using index notation, this becomes:

\begin{equation}
\boxed{
\Delta_T\phi
=
\nabla^{T}_{a}\nabla^{T\,a}\phi .
}
\end{equation}

Indices are raised using the regularized inverse metric wherever the ordinary inverse fails at the TZP.

\paragraph{Interpretation.}
$\Delta_T$ is the Laplacian adapted to TZP geometry. It measures second-order field variation using the Trawinistic connection.

\par\vspace{0.35em}
\noindent\textbf{Derivable Consequences.}\quad
\begingroup
\footnotesize
\raggedright
The canonical equation anchors Constrained Dipole Field Configuration by connecting the source condition to the derivation layer and proof certificate.
\par\endgroup

\tzpsection{Equation Support}
\noindent\textbf{1. Dipole Locations}\par
\[
\mathbf{r}_{+}=a\hat{\mathbf{z}},\qquad \mathbf{r}_{-}=-a\hat{\mathbf{z}}
\]
\noindent This fixes the antipodal locations of the constrained dipole pair along the carrier axis.\par
\noindent\textbf{2. Single Dipole Field}\par
\[
\mathbf{B}_{\mathrm{dip}}(\mathbf{r};\mathbf{m})=\frac{\mu_{0}}{4\pi}\frac{3\hat{\mathbf{r}}(\mathbf{m}\cdot\hat{\mathbf{r}})-\mathbf{m}}{\lVert\mathbf{r}\rVert^{3}}
\]
\noindent This states the Biot--Savart dipole field contributed by a single localized moment.\par
\noindent\textbf{3. Total Constrained Field}\par
\[
\mathbf{B}=\mathbf{B}_{\mathrm{dip}}(\mathbf{r}-\mathbf{r}_{+};\mathbf{m})+\mathbf{B}_{\mathrm{dip}}(\mathbf{r}-\mathbf{r}_{-};-\mathbf{m})
\]
\noindent This closes the entry by assembling the total constrained field from both dipoles and the boundary constraint sector.\par

\clearpage
\noindent\textbf{[ID 0318] Constrained Dipole Field Configuration}\par
\vspace{0.35em}
\tzpsection{Dictionary}
\begingroup\small
\tzpsection{Definition}
Constrained Dipole Field Configuration is a registry definition in Field Theory and Action Principles with secondary pressure from Manifold and Metric Dynamics.

% BEGIN TZP CONTEXTUAL MEANING
\tzpsection{Contextual Meaning}
\begingroup\footnotesize\setlength{\emergencystretch}{3em}\sloppy
\textbf{Formal statement:} Two magnetic dipoles of equal magnitude mu sub 0 , oriented along the z-axis with opposite directions, are located at r sub 1 =(0,0,-d/2), ; r sub 2 =(0,0,d/2). The field produced by a single dipole at position r sub 0 is B sub dip (r;mu,r sub 0 ) = ratio mu sub 0 4pi [ ratio 3(r-r sub 0 )bigl(mudot(r-r sub 0 )bigr) lvert r-r sub 0 rvert to the power 5 - ratio mu lvert r-r sub 0 rvert to the power 3 ]. The total internal field is B sub total (r,t) = B sub dip (r;mu sub 1 (t),r sub 1 ) + B sub dip (r;mu sub 2 (t),r sub 2 ) + B sub constraint (r,t). \textbf{Contextual meaning:} The entry reads Constrained Dipole Field Configuration as a controlled technical headword whose equation layer, proof route, and registry role are interpreted together. \textbf{Derivable consequences:} The entry now states the constrained dipole configuration as a real field construction instead of truncated prose fragments. \textbf{Lemma support:} Dipole Locations; Single Dipole Field; Total Constrained Field.\par
\endgroup
% END TZP CONTEXTUAL MEANING

\tzpsection{Technical Interpretation}
sub T is the Laplacian adapted to TZP geometry. It measures second-order field variation using the Trawinistic connection.

\tzpsection{Equation Grounding}
Equation anchors: target = C:azlan-assistanthandoffinboxai7.gpt5 sub research sub fs sub watcher.json ; target.tmp findings :[ new content ], recommended sub tasks :[] | Set-Content; and mathbf v sub in ( x sub 0)= - v sub out ( x sub 0),. Proof route: show that the stated dipole locations and single-dipole field law combine into the constrained total field configuration. Formalization route: prove that superposition of the two dipole sectors and the constraint sector yields the stated total field.

\tzpsection{Interpretive Role}
The entry now states the constrained dipole configuration as a real field construction instead of truncated prose fragments.
\endgroup

\tzpsection{TZP Thesaurus}
\begin{enumerate}[leftmargin=1.6em,itemsep=6pt,topsep=4pt]
\item \textbf{Antipodal Magnetic Moments}: The spatial arrangement of two equal and opposite magnetic point sources locked to a fixed geometric axis.
\item \textbf{Biot-Savart Superposition}: The summation of individual vector magnetic fields to produce a single, unified flux density map.
\item \textbf{Constrained Axial Magnetostatics}: The imposition of strict boundary conditions ensuring field lines remain trapped within a localized geometric channel.
\end{enumerate}

\tzpsection{Encyclopedia Mathematica}
\begingroup\footnotesize\sloppy
The visual plate below is reserved for the source-specific equation and progression graphic for [ID 0318]. It is included only as a companion to the canonical equation, not as a substitute for the formal text.
\par\endgroup
\ifTZPDarkEdition
\tzpdictionaryvisualband{D:/TZPID/kdp_workflow/build/tikz_figures/ID0318_dictionary_figure_dark.pdf}
\fi

\clearpage
\begingroup\tzpsupportsetup
\noindent\textbf{\fontsize{10.8pt}{12.8pt}\selectfont [ID 0318] Constructive Logic and Proof Certificate}\par
\noindent\textbf{\fontsize{10pt}{12pt}\selectfont Constrained Dipole Field Configuration}\par
\vspace{0.45em}
\begingroup
\footnotesize
\setlength{\parskip}{0.22em}
\setlength{\parindent}{0pt}
\noindent\textbf{Curry--Howard--Lambek Interpretation}\par
Under the Curry--Howard--Lambek correspondence, this registry entry is read as a typed proof
object: the stated source condition supplies the proposition, the derivation supplies the
morphism, and the proof certificate records the machine handles needed to check the obligation
in the formal registry.

\vspace{0.45em}
\noindent\textbf{CHL Proof}\par
\textbf{Statement:} two magnetic dipoles of equal magnitude \_0, oriented along the z-axis with opposite
directions, are located at The field produced by a single d...

\textbf{Logic Validation:}\\
\textbf{Proposition:} ``Constrained Dipole Field Configuration is valid as a formal registry statement when its
source equation and semantic claim are read together.''\\
\textbf{Proof:} The source semantics state that two magnetic dipoles of equal magnitude \_0, oriented
along the z-axis with opposite directions, are located at The field produced by a
single dipole at position r\_0 is The total internal field is. \_T is the Laplacian
adapted to TZP geometry. It measures second-order field variation using the Trawinistic
connection.. The CHL validation treats the equation as the formal anchor and the
semantic text as the interpretation constraint.

\textbf{VERDICT:} \ensuremath{\checkmark}\ \textbf{TRUE} --- source-backed formal validation

\textbf{Formal Status:} \ensuremath{\checkmark}\ THEORETICAL -- source-backed CHL proof obligation

\textbf{Physical Status:} THEORETICAL -- requires model-specific empirical interpretation
\par\vspace{0.45em}
\noindent{\tiny\textbf{Isabelle/HOL.} \tzpplaincode{physics\_validity\_from\_model, external\_validity\_package, registry\_number\_matches, equation\_count\_matches}}\par
\ifTZPDarkEdition
\tzpproofvisualband{D:/TZPID/kdp_workflow/build/tikz_figures/ID0318_proof_certificate_figure_dark.pdf}
\fi
\endgroup
\endgroup
